% Obsahový zdroj: import_hodiny/38_Pascaluv_trojuhelnik.tex a odpovídající PDF (obě strany).
% P1–P5: zdrojové Příklady 1–5; C1–C5: zdrojové Procvičování 1–5.
% Matematický obsah zachován; výklad, odvození a řešení pedagogicky rozvedeny.
% Upřesnění a opravy: viz KMB_29-41_Kontrola.md; původní import zůstává beze změny.
% Rozpis: 4_rocnik/tematicke_plany/Rozpis_vyuky_MAT_4_rocnik.tex, hodina 38.
% Plán: 4_rocnik/tematicke_plany/SPS.pdf, oddíly 2.1–2.3.
% Sazba: sablony/pripravy/TEMPLATE_Priprava.tex a společný priprava.sty.
% Vlastní doplnění: didaktický výklad, zdůvodnění, metodika; žádné nové číslované úlohy.
% Učebnice/09_Kombinatorika nemá obsahovou kapitolu; místní kombinatorické listy ani Realisticky.cz nejsou dostupné.
\documentclass[11pt,a4paper]{article}
% Spolecny styl pro ucitelske pripravy MAT (Work/MAT)
% Verze 1.1, 2026-10-23
\NeedsTeXFormat{LaTeX2e}
\ProvidesPackage{priprava}[2026/10/23 v1.1 Sablona ucitelskych priprav MAT]

% Zakladni typografie a cestina
\RequirePackage[T1]{fontenc}
\RequirePackage[utf8]{inputenc}
\RequirePackage[czech]{babel}
\RequirePackage{lmodern}
\RequirePackage{microtype}
\RequirePackage{amsmath,amssymb}
\RequirePackage{xcolor}
\RequirePackage{enumitem}
\RequirePackage[most]{tcolorbox}
\RequirePackage{tabularx,booktabs,array}
\RequirePackage{fancyhdr}
\RequirePackage{needspace}
\RequirePackage{geometry}
\RequirePackage[hidelinks]{hyperref}

\geometry{a4paper,top=17mm,bottom=18mm,left=23mm,right=23mm}
\definecolor{notegray}{gray}{0.39}
\setlength{\parindent}{0pt}
\setlength{\parskip}{2.5pt plus 1pt minus 0.5pt}
\setlist[itemize]{topsep=2pt,itemsep=1pt,parsep=0pt,leftmargin=1.45em}
\setlist[enumerate]{topsep=2pt,itemsep=2pt,parsep=0pt,leftmargin=1.65em}
\setcounter{secnumdepth}{0}
\renewcommand{\arraystretch}{1.12}

% Zahlavi: vlevo rocnik | globalni tema, vpravo cislo: nazev hodiny
% Zapati zustava zcela prazdne.
\pagestyle{fancy}
\fancyhf{}
\renewcommand{\headrulewidth}{0.35pt}
\renewcommand{\footrulewidth}{0pt}
\setlength{\headheight}{14pt}
\newcommand{\nastavPripravu}[4]{%
  \fancyhead[L]{\small #1\enspace\textbar\enspace #2}%
  \fancyhead[R]{\small #3: #4}%
}

% Minimalisticke cernobile boxy. Kazda definice zvlast.
\tcbset{commonlesson/.style={enhanced,breakable=false,colback=white,colframe=black,
  colbacktitle=white,coltitle=black,boxrule=0.45pt,
  arc=0.8mm,outer arc=0.8mm,left=3mm,right=3mm,top=1.6mm,bottom=1.6mm,
  toptitle=1mm,bottomtitle=0.7mm,fonttitle=\bfseries,
  attach boxed title to top left={xshift=2.5mm,yshift=-1.15mm},
  boxed title style={colback=white,colframe=white,boxrule=0pt,left=1mm,right=1mm}}}
\newtcolorbox{zadaniBox}[1][]{commonlesson,title={Zadání},#1}
\newtcolorbox{prikladBox}[1][]{commonlesson,title={Příklad},#1}
\newtcolorbox{definiceBox}[1][]{commonlesson,title={Definice},#1}

% Spolecny box pro zadani a jeho reseni, bez mezery mezi castmi.
% Pouziti:
%   \begin{prikladReseniBox}
%     Zadani ulohy.
%     \tcblower
%     \begin{align*} ... \end{align*}
%   \end{prikladReseniBox}
% Spodni cast ma jemny sedy podklad a delici linku; lze v ni sazet
% vice radku matematiky. Obe casti maji jen jeden spolecny ramecek.
\tcbset{solutionlesson/.style={
  commonlesson,
  enhanced jigsaw,
  breakable,
  colbacklower=black!3,
  segmentation style={solid,draw=black!25,line width=0.35pt},
  fontlower=\small\color{notegray},
  before lower={\noindent\textbf{Řešení:}\enspace},
  leftlower=3mm,rightlower=3mm
}}
\newtcolorbox{prikladReseniBox}[1][]{solutionlesson,title={Příklad},#1}
\newtcolorbox{zadaniReseniBox}[1][]{solutionlesson,title={Zadání},#1}

% Pomucky pro autora pripravy: sedive poznamky; \reseni pro starsi pripravy
\newcommand{\poznamka}[1]{\par\begingroup\small\color{notegray}\itshape
  \textnormal{\textbf{Poznámka pro učitele.}}\enspace #1\par\endgroup}
\newcommand{\reseni}[1]{\par\begingroup\small\color{notegray}
  \textbf{Očekávané řešení:}\enspace #1\par\endgroup}
\newcommand{\blok}[1]{\needspace{5\baselineskip}\subsection*{#1}}
\newcommand{\vyrok}[1]{\uv{#1}}

\nastavPripravu{4. ročník}{Kombinatorika}{38}{Pascalův trojúhelník}
\begin{document}


Připomeň Pascalův vztah z hodiny 37 a nech z krajních jedniček postupně sestavit další řádky. Zápis zviditelní symetrii a umožní výpočet bez opakovaného dosazování faktoriálů.
\poznamka{Předpoklady: Pascalův vztah, symetrie, součet řádku, kombinace a permutace s opakováním. Cíl: žák správně indexuje trojúhelník, vysvětlí jeho konstrukci a počítá monotónní cesty i cesty přes pevný bod. Důraz je na propojení tří pohledů: součty, výběr pozic kroků, mřížka.}

\blok{Výklad a odvození}
\begin{definiceBox}
\textbf{Pascalův trojúhelník} je uspořádání kombinačních čísel do řádků. Řádek s indexem $n\in\mathbb N_0$ obsahuje postupně $\binom n0,\binom n1,\ldots,\binom nn$. Má $n+1$ členů; první řádek má index $0$.
\end{definiceBox}

Na krajích je vždy jedna; každý vnitřní člen je součtem dvou členů bezprostředně nad ním. Platnost konstrukce zaručuje Pascalův vztah, krajní hodnoty zaručuje $\binom n0=\binom nn=1$.
\begin{center}
\begin{tikzpicture}
 \node at (-2.2,0.5) {$n$};
 \node at (0,0.5) {čísla v řádku};
 \draw (-1.85,0.75)--(-1.85,-2.9);
 \foreach \n/\hodnoty in {0/{1},1/{1,1},2/{1,2,1},3/{1,3,3,1},4/{1,4,6,4,1},5/{1,5,10,10,5,1}} {
   \node at (-2.2,-0.52*\n) {$\n$};
   \foreach \v [count=\k from 0] in \hodnoty {
     \node at (0.62*\k-0.31*\n,-0.52*\n) {$\v$};
   }
 }
\end{tikzpicture}
\end{center}
Zrcadlová shoda čísel odpovídá symetrii; součet řádku $n$ je $2^n$. Index $k$ počítáme zleva od nuly, nikoli od jedné.\begin{definiceBox}
\textbf{Monotónní cesta v jednotkové mřížce} zde používá pouze jednotkové kroky doprava $D$ a nahoru $N$. Cesta z $(0,0)$ do $(a,b)$ pro celá $a,b\geq0$ má právě $a+b$ kroků a je nejkratší cestou po hranách mřížky mezi těmito body.
\end{definiceBox}

Každá cesta je řetězec s $a$ znaky D a $b$ znaky N. Vybereme pozice kroků D a ostatní jsou určeny. Proto
\[
 \text{počet cest}=\binom{a+b}{a}=\binom{a+b}{b}
 =\frac{(a+b)!}{a!\,b!}.
\]
Tím se výběr pozic shoduje s permutací opakovaných kroků. I mřížka sama dává Pascalův vztah: do vnitřního bodu přijdeme buď zleva, nebo zdola. Tyto dvě skupiny cest jsou disjunktní a jejich počty se sčítají. Na osách je do každého bodu pouze jedna monotónní cesta.

Má-li cesta projít bodem $(u,v)$, musí být $0\leq u\leq a$, $0\leq v\leq b$. Volbu cesty k tomuto bodu a od něj násobíme, protože každá dvojice úseků určí právě jednu celou cestu. Pro bod mimo tento obdélník není taková monotónní cesta možná.


\poznamka{Průběh: 0–10 min sestavení řádků a P1; 10–18 min indexy, P2 a P5; 18–31 min P3–P4; 31–42 min C1–C5; 42–45 min kontrola. P5 je krátké čtení z trojúhelníku, nikoli další dlouhý výpočet. Při pomalejším tempu přesuň C5 do domácí práce.}

\blok{Společně řešené příklady}

% Zdroj P1: import_hodiny/38_Pascaluv_trojuhelnik.tex, Příklad 1.
\begin{prikladReseniBox}[title={Příklad P1}]
Doplň řádek \(n=6\) Pascalova trojúhelníku.
\tcblower
Vyjdeme z řádku $n=5$: $1,5,10,10,5,1$. Doplníme krajní jedničky a mezi ně součty sousedních hodnot: $1+5=6$, $5+10=15$, $10+10=20$, $10+5=15$, $5+1=6$. Výsledný řádek $n=6$ je $1,6,15,20,15,6,1$. Má sedm členů, je symetrický a jeho součet je $64=2^6$.
\end{prikladReseniBox}


% Zdroj P2: import_hodiny/38_Pascaluv_trojuhelnik.tex, Příklad 2.
\begin{prikladReseniBox}[title={Příklad P2}]
Urči součet čísel v řádku \(n=5\).
\tcblower
Součet \(1+5+10+10+5+1=32\). Také přímo \(2^5=32\), protože řádek tvoří \(\binom5k\).
\end{prikladReseniBox}


% Zdroj P3: import_hodiny/38_Pascaluv_trojuhelnik.tex, Příklad 3.
\begin{prikladReseniBox}[title={Příklad P3}]
Kolik je nejkratších cest z \((0,0)\) do \((3,4)\), pouze doprava a nahoru?
\tcblower
Každá cesta je posloupnost sedmi kroků, v níž jsou právě tři \(D\) a čtyři \(N\). Výběr tří pozic pro \(D\): \(\binom73=35\).
\end{prikladReseniBox}


% Zdroj P4: import_hodiny/38_Pascaluv_trojuhelnik.tex, Příklad 4.
\begin{prikladReseniBox}[title={Příklad P4}]
Kolik monotónních cest z $(0,0)$ do $(4,3)$ prochází bodem $(2,1)$?
\tcblower
Od $(0,0)$ k $(2,1)$ potřebujeme dvě D a jedno N: $\binom32=3$ cesty. Od $(2,1)$ k $(4,3)$ se souřadnice změní o $(2,2)$: potřebujeme dvě D a dvě N, tedy $\binom42=6$ cest. Každý první úsek lze spojit s každým druhým, takže celkem $3\cdot6=18$. Průchod bodem jednoznačně rozděluje celou cestu, proto žádnou nezapočítáme dvakrát.
\end{prikladReseniBox}

\begin{center}
\begin{tikzpicture}[scale=0.62]
 \draw[step=1,black!25] (0,0) grid (4,3);
 \draw[->] (0,0)--(4.7,0) node[right] {$x$};
 \draw[->] (0,0)--(0,3.6) node[above] {$y$};
 \draw[very thick,->] (0,0)--(1,0)--(1,1)--(2,1);
 \draw[very thick,->] (2,1)--(3,1)--(3,2)--(4,2)--(4,3);
 \fill (0,0) circle (2pt) node[below left] {$O(0,0)$};
 \fill (2,1) circle (2pt) node[below right] {$P(2,1)$};
 \fill (4,3) circle (2pt) node[above right] {$Q(4,3)$};
\end{tikzpicture}
\end{center}
\poznamka{Náčrt ukazuje jednu přípustnou cestu k P a její pokračování k Q. Nejde o výpis všech cest; číslo 18 získáme násobením počtů obou úseků.}

% Zdroj P5: import_hodiny/38_Pascaluv_trojuhelnik.tex, Příklad 5.
\begin{prikladReseniBox}[title={Příklad P5}]
Najdi v řádku \(n=8\) číslo se spodním indexem \(k=2\).
\tcblower
Je to \(\binom82=28\). Můžeme ho získat Pascalovým součtem \(\binom71+\binom72=7+21=28\).
\end{prikladReseniBox}


\newpage
\blok{Samostatná práce a zásoba k procvičení}

Řádky vždy označ indexem $n$. U C4–C5 i všech cest v této přípravě jsou povoleny pouze jednotkové kroky doprava a nahoru po hranách mřížky.

% Zdroj C1: import_hodiny/38_Pascaluv_trojuhelnik.tex, Procvičování, položka 1.
\begin{zadaniReseniBox}[title={Úloha C1}]
Vypiš řádek \(n=4\).
\tcblower
Z řádku $1,3,3,1$ vytvoříme součty $1+3,3+3,3+1$ a přidáme kraje: řádek $n=4$ je $1,4,6,4,1$.
\end{zadaniReseniBox}

% Zdroj C2: import_hodiny/38_Pascaluv_trojuhelnik.tex, Procvičování, položka 2.
\begin{zadaniReseniBox}[title={Úloha C2}]
Urči střední číslo řádku \(n=6\).
\tcblower
Řádek $n=6$ má sedm členů. Střední je čtvrtý zleva, tedy s indexem $k=3$: $\binom63=20$.
\end{zadaniReseniBox}

% Zdroj C3: import_hodiny/38_Pascaluv_trojuhelnik.tex, Procvičování, položka 3.
\begin{zadaniReseniBox}[title={Úloha C3}]
Urči součet řádku \(n=7\).
\tcblower
Součet celého řádku je podle počtu všech podmnožin $2^7=128$.
\end{zadaniReseniBox}

% Zdroj C4: import_hodiny/38_Pascaluv_trojuhelnik.tex, Procvičování, položka 4.
\begin{zadaniReseniBox}[title={Úloha C4}]
Kolik monotónních nejkratších cest vede z $(0,0)$ do $(2,3)$, povolujeme-li jen jednotkové kroky doprava a nahoru?
\tcblower
Potřebujeme dvě D a tři N, celkem pět kroků. Pozice dvou D vybíráme $\binom52=10$ způsoby; ostatní pozice obsadí N.
\end{zadaniReseniBox}

% Zdroj C5: import_hodiny/38_Pascaluv_trojuhelnik.tex, Procvičování, položka 5.
\begin{zadaniReseniBox}[title={Úloha C5}]
Kolik monotónních nejkratších cest vede z $(0,0)$ do $(4,2)$ přes $(2,1)$, povolujeme-li jen jednotkové kroky doprava a nahoru?
\tcblower
K bodu $(2,1)$ vedou tři kroky, z toho dvě D: $\binom32=3$ možnosti. Druhý úsek má opět posun $(2,1)$, tedy další tři možnosti. Počty úseků násobíme: $3\cdot3=9$.
\end{zadaniReseniBox}


\poznamka{Typické chyby: zaměnění indexu řádku a jeho pořadového čísla; zaměnění bodů za kroky; sčítání počtů dvou po sobě jdoucích úseků. U P4 je sedm kroků celkem, ale rozdělují se na tři a čtyři. Příště se řádky trojúhelníku objeví jako koeficienty rozvoje mocnin dvojčlenu.}
\end{document}
